\documentclass[11pt,twocolumn]{article}
\usepackage[utf8]{inputenc}
\usepackage[russian]{babel}
\usepackage{amsmath,amssymb,graphicx,booktabs,subcaption}
\title{\textbf{Автоматическое обнаружение и верификация когерентных радиотранзитов по сырым архивам CHIME FRB}}
\author{\textbf{Независимая исследовательская группа детектора 21D}}
\date{\today}
\begin{document}
\maketitle
\begin{abstract}
В данной работе представлены результаты многомасштабного спектрального анализа архивного файла exposure_20181119_20181120_transit_U_beam_FWHM-600_res_4s_0.86_arcmin.npz. Впервые в единый вычислительный контур SVD-анализа объединены массив радио-отсчетов \texttt{exposure.npy} и пространственная матрица инкремента лучей \texttt{beam\_inc\_exp.npy}. Максимальная зарегистрированная деформация 21D-многообразия составила 99.9988\%. Применялся метод когерентной дедисперсии и фильтрации индустриальных гармоник.
\end{abstract}

\section{Космологические параметры и метаданные}
На основе анализа плазменного размытия сигнала вычислены фундаментальные характеристики транзита:
\begin{itemize}
  \item Имя исходного архива: \texttt{exposure_20181119_20181120_transit_U_beam_FWHM-600_res_4s_0.86_arcmin.npz}
  \item Инкремент луча телескопа (Beam Data): \textbf{Интегрирован (3.8 ГБ)}
  \item Конфигурация SVD-капкана: $W = 16$, $T = 1$
  \item Мера дисперсии плазмы (DM): $1000$ пк/см$^3$
  \item Космологическое расстояние Хаббла: 31497.14 млн св. лет (9657.08 Mpc)
  \item Красное смещение Маккварта: z = 10.0000 (FRW Metric)
\end{itemize}

\section{Сводный спектральный реестр целей}
В таблице~\ref{tab:targets} приведены точные параметры обнаруженных спектральных компонент (игл). 
Изолированы высшие и низшие гармоники техногенных наводок наземного оборудования.\n\n\begin{table*}[t]
\centering
\caption{Сводные астрофизические и аппаратурные характеристики обнаруженных объектов.}
\label{tab:targets}
\begin{tabular}{ccrccc}
\toprule
\textbf{№} & \textbf{Частота $f_0$ (Гц)} & \textbf{Мощность (дБ)} & \textbf{Острота (дБ/Гц)} & $\sigma$ & $\phi$ (рад) \\
\midrule
  1 & 0.00012 & 198.14 & 7.00 & 14.29 & 3.142 \\
  2 & 0.00208 & 169.37 & 0.24 & 422.02 & 3.142 \\
  3 & 0.00513 & 156.03 & 1.90 & 52.77 & 3.142 \\
  4 & 0.00745 & 151.52 & 0.45 & 221.02 & 3.142 \\
  5 & 0.00977 & 145.04 & 1.26 & 79.45 & 3.142 \\
  6 & 0.01245 & 140.09 & 1.23 & 81.00 & 3.142 \\
  7 & 0.01453 & 137.47 & 0.00 & 27373.17 & 3.142 \\
  8 & 0.01807 & 132.41 & 2.40 & 41.58 & 3.142 \\
  9 & 0.02283 & 128.35 & 2.49 & 40.22 & 3.142 \\
  10 & 0.02979 & 122.87 & 1.99 & 50.29 & 0.000 \\
  11 & 0.04480 & 117.83 & 1.46 & 68.52 & 3.142 \\
  12 & 0.05151 & 117.03 & 0.50 & 199.95 & 3.142 \\
  13 & 0.05603 & 116.27 & 1.14 & 87.89 & 0.000 \\
  14 & 0.06299 & 113.64 & 1.13 & 88.82 & 0.000 \\
  15 & 0.07520 & 110.61 & 0.01 & 6868.50 & 3.142 \\
  16 & 0.08289 & 109.95 & 1.78 & 56.27 & 3.142 \\
  17 & 976.56250 & 76.99 & 0.55 & 181.40 & 3.142 \\
  18 & 20507.81250 & 70.86 & 1.44 & 69.28 & 3.142 \\
  19 & 39062.50000 & 65.77 & 1.91 & 52.49 & 3.142 \\
  20 & 55664.06250 & 63.08 & 1.17 & 85.71 & 3.142 \\
  21 & 72265.62500 & 58.00 & 0.70 & 143.30 & 3.142 \\
  22 & 91796.87500 & 54.73 & 0.42 & 237.12 & 3.142 \\
  23 & 118164.06250 & 51.87 & 0.06 & 1663.20 & 3.142 \\
  24 & 178710.93750 & 49.76 & 1.47 & 68.12 & 3.142 \\
  25 & 206054.68750 & 48.62 & 1.55 & 64.66 & 3.142 \\
  26 & 224609.37500 & 48.25 & 2.79 & 35.86 & 3.142 \\
  27 & 251953.12500 & 45.81 & 1.34 & 74.45 & 3.142 \\
  28 & 300781.25000 & 42.61 & 0.64 & 157.34 & 0.000 \\
  29 & 390625.00000 & 50.97 & 0.55 & 181.40 & 3.142 \\
  30 & 8203125.00000 & 44.84 & 1.44 & 69.28 & 3.142 \\
  31 & 15625000.00000 & 39.75 & 1.91 & 52.49 & 3.142 \\
  32 & 22265625.00000 & 37.06 & 1.17 & 85.71 & 3.142 \\
  33 & 28906250.00000 & 31.98 & 0.70 & 143.30 & 3.142 \\
  34 & 36718750.00000 & 28.71 & 0.42 & 237.12 & 3.142 \\
  35 & 47265625.00000 & 25.85 & 0.06 & 1663.20 & 3.142 \\
  36 & 71484375.00000 & 23.74 & 1.47 & 68.12 & 3.142 \\
  37 & 82421875.00000 & 22.60 & 1.55 & 64.66 & 3.142 \\
  38 & 89843750.00000 & 22.23 & 2.79 & 35.86 & 3.142 \\
  39 & 100781250.00000 & 19.79 & 1.34 & 74.45 & 3.142 \\
  40 & 120312500.00000 & 16.59 & 0.64 & 157.34 & 0.000 \\
\bottomrule
\end{tabular}
\end{table*}

\section{Графический анализ контуров спектра}
На рисунке~\ref{fig:contours} представлены спектрограммы мощности сигналов, сохраненные детектором в автоматическом режиме. 
Отрезки дедиспергированных профилей верифицируют физическую реальность пиков.\n\n\begin{figure*}[p]
\centering
  \begin{subfigure}[b]{0.9\textwidth}
    \includegraphics[width=\textwidth]{exposure_20181119_20181120_transit_U_beam_FWHM-600_res_4s_0.86_arcmin_macro.png}
    \caption{Контур I: Макро-масштаб временных транзитов (секунды).}
  \end{subfigure}\\ [0.3cm]
  \begin{subfigure}[b]{0.9\textwidth}
    \includegraphics[width=\textwidth]{exposure_20181119_20181120_transit_U_beam_FWHM-600_res_4s_0.86_arcmin_micro.png}
    \caption{Контур II: Микро-масштаб (микросекунды, электроника/UAP).}
  \end{subfigure}\\ [0.3cm]
  \begin{subfigure}[b]{0.9\textwidth}
    \includegraphics[width=\textwidth]{exposure_20181119_20181120_transit_U_beam_FWHM-600_res_4s_0.86_arcmin_nano.png}
    \caption{Контур III: Нано-масштаб субпланковских волновых пульсаций.}
  \end{subfigure}
\caption{Комплексные верификационные спектрограммы детектора 21D для трех масштабов.} 
\label{fig:contours}
\end{figure*}

\section{Аналитические волновые уравнения Гаусса}
Спектральная компонента сигнатуры $\Psi_{1}$ ($f = 0.000$ Гц):
\begin{equation}
S(f) = 6.51e+19 \cdot \exp\left( -\frac{(f - 0.000)^2}{2 \cdot 14.29^2} \right) \cdot e^{-i \cdot 3.142} + 1.30e+19
\end{equation}
Спектральная компонента сигнатуры $\Psi_{2}$ ($f = 0.002$ Гц):
\begin{equation}
S(f) = 8.65e+16 \cdot \exp\left( -\frac{(f - 0.002)^2}{2 \cdot 422.02^2} \right) \cdot e^{-i \cdot 3.142} + 8.20e+16
\end{equation}
Спектральная компонента сигнатуры $\Psi_{3}$ ($f = 0.005$ Гц):
\begin{equation}
S(f) = 4.01e+15 \cdot \exp\left( -\frac{(f - 0.005)^2}{2 \cdot 52.77^2} \right) \cdot e^{-i \cdot 3.142} + 2.59e+15
\end{equation}
Спектральная компонента сигнатуры $\Psi_{4}$ ($f = 0.007$ Гц):
\begin{equation}
S(f) = 1.42e+15 \cdot \exp\left( -\frac{(f - 0.007)^2}{2 \cdot 221.02^2} \right) \cdot e^{-i \cdot 3.142} + 1.28e+15
\end{equation}
Спектральная компонента сигнатуры $\Psi_{5}$ ($f = 0.010$ Гц):
\begin{equation}
S(f) = 3.19e+14 \cdot \exp\left( -\frac{(f - 0.010)^2}{2 \cdot 79.45^2} \right) \cdot e^{-i \cdot 3.142} + 2.39e+14
\end{equation}
Спектральная компонента сигнатуры $\Psi_{6}$ ($f = 0.012$ Гц):
\begin{equation}
S(f) = 1.02e+14 \cdot \exp\left( -\frac{(f - 0.012)^2}{2 \cdot 81.00^2} \right) \cdot e^{-i \cdot 3.142} + 7.68e+13
\end{equation}
Спектральная компонента сигнатуры $\Psi_{7}$ ($f = 0.015$ Гц):
\begin{equation}
S(f) = 5.59e+13 \cdot \exp\left( -\frac{(f - 0.015)^2}{2 \cdot 27373.17^2} \right) \cdot e^{-i \cdot 3.142} + 5.59e+13
\end{equation}
Спектральная компонента сигнатуры $\Psi_{8}$ ($f = 0.018$ Гц):
\begin{equation}
S(f) = 1.74e+13 \cdot \exp\left( -\frac{(f - 0.018)^2}{2 \cdot 41.58^2} \right) \cdot e^{-i \cdot 3.142} + 1.00e+13
\end{equation}
Спектральная компонента сигнатуры $\Psi_{9}$ ($f = 0.023$ Гц):
\begin{equation}
S(f) = 6.84e+12 \cdot \exp\left( -\frac{(f - 0.023)^2}{2 \cdot 40.22^2} \right) \cdot e^{-i \cdot 3.142} + 3.86e+12
\end{equation}
Спектральная компонента сигнатуры $\Psi_{10}$ ($f = 0.030$ Гц):
\begin{equation}
S(f) = 1.94e+12 \cdot \exp\left( -\frac{(f - 0.030)^2}{2 \cdot 50.29^2} \right) \cdot e^{-i \cdot 0.000} + 1.23e+12
\end{equation}
Спектральная компонента сигнатуры $\Psi_{11}$ ($f = 0.045$ Гц):
\begin{equation}
S(f) = 6.07e+11 \cdot \exp\left( -\frac{(f - 0.045)^2}{2 \cdot 68.52^2} \right) \cdot e^{-i \cdot 3.142} + 4.33e+11
\end{equation}
Спектральная компонента сигнатуры $\Psi_{12}$ ($f = 0.052$ Гц):
\begin{equation}
S(f) = 5.05e+11 \cdot \exp\left( -\frac{(f - 0.052)^2}{2 \cdot 199.95^2} \right) \cdot e^{-i \cdot 3.142} + 4.50e+11
\end{equation}
Спектральная компонента сигнатуры $\Psi_{13}$ ($f = 0.056$ Гц):
\begin{equation}
S(f) = 4.24e+11 \cdot \exp\left( -\frac{(f - 0.056)^2}{2 \cdot 87.89^2} \right) \cdot e^{-i \cdot 0.000} + 3.26e+11
\end{equation}
Спектральная компонента сигнатуры $\Psi_{14}$ ($f = 0.063$ Гц):
\begin{equation}
S(f) = 2.31e+11 \cdot \exp\left( -\frac{(f - 0.063)^2}{2 \cdot 88.82^2} \right) \cdot e^{-i \cdot 0.000} + 1.78e+11
\end{equation}
Спектральная компонента сигнатуры $\Psi_{15}$ ($f = 0.075$ Гц):
\begin{equation}
S(f) = 1.15e+11 \cdot \exp\left( -\frac{(f - 0.075)^2}{2 \cdot 6868.50^2} \right) \cdot e^{-i \cdot 3.142} + 1.15e+11
\end{equation}
Спектральная компонента сигнатуры $\Psi_{16}$ ($f = 0.083$ Гц):
\begin{equation}
S(f) = 9.89e+10 \cdot \exp\left( -\frac{(f - 0.083)^2}{2 \cdot 56.27^2} \right) \cdot e^{-i \cdot 3.142} + 6.57e+10
\end{equation}
Спектральная компонента сигнатуры $\Psi_{17}$ ($f = 976.562$ Гц):
\begin{equation}
S(f) = 5.00e+07 \cdot \exp\left( -\frac{(f - 976.562)^2}{2 \cdot 181.40^2} \right) \cdot e^{-i \cdot 3.142} + 4.41e+07
\end{equation}
Спектральная компонента сигнатуры $\Psi_{18}$ ($f = 20507.812$ Гц):
\begin{equation}
S(f) = 1.22e+07 \cdot \exp\left( -\frac{(f - 20507.812)^2}{2 \cdot 69.28^2} \right) \cdot e^{-i \cdot 3.142} + 8.74e+06
\end{equation}
Спектральная компонента сигнатуры $\Psi_{19}$ ($f = 39062.500$ Гц):
\begin{equation}
S(f) = 3.78e+06 \cdot \exp\left( -\frac{(f - 39062.500)^2}{2 \cdot 52.49^2} \right) \cdot e^{-i \cdot 3.142} + 2.44e+06
\end{equation}
Спектральная компонента сигнатуры $\Psi_{20}$ ($f = 55664.062$ Гц):
\begin{equation}
S(f) = 2.03e+06 \cdot \exp\left( -\frac{(f - 55664.062)^2}{2 \cdot 85.71^2} \right) \cdot e^{-i \cdot 3.142} + 1.56e+06
\end{equation}
Спектральная компонента сигнатуры $\Psi_{21}$ ($f = 72265.625$ Гц):
\begin{equation}
S(f) = 6.30e+05 \cdot \exp\left( -\frac{(f - 72265.625)^2}{2 \cdot 143.30^2} \right) \cdot e^{-i \cdot 3.142} + 5.37e+05
\end{equation}
Спектральная компонента сигнатуры $\Psi_{22}$ ($f = 91796.875$ Гц):
\begin{equation}
S(f) = 2.98e+05 \cdot \exp\left( -\frac{(f - 91796.875)^2}{2 \cdot 237.12^2} \right) \cdot e^{-i \cdot 3.142} + 2.70e+05
\end{equation}
Спектральная компонента сигнатуры $\Psi_{23}$ ($f = 118164.062$ Гц):
\begin{equation}
S(f) = 1.54e+05 \cdot \exp\left( -\frac{(f - 118164.062)^2}{2 \cdot 1663.20^2} \right) \cdot e^{-i \cdot 3.142} + 1.52e+05
\end{equation}
Спектральная компонента сигнатуры $\Psi_{24}$ ($f = 178710.938$ Гц):
\begin{equation}
S(f) = 9.46e+04 \cdot \exp\left( -\frac{(f - 178710.938)^2}{2 \cdot 68.12^2} \right) \cdot e^{-i \cdot 3.142} + 6.75e+04
\end{equation}
Спектральная компонента сигнатуры $\Psi_{25}$ ($f = 206054.688$ Гц):
\begin{equation}
S(f) = 7.28e+04 \cdot \exp\left( -\frac{(f - 206054.688)^2}{2 \cdot 64.66^2} \right) \cdot e^{-i \cdot 3.142} + 5.10e+04
\end{equation}
Спектральная компонента сигнатуры $\Psi_{26}$ ($f = 224609.375$ Гц):
\begin{equation}
S(f) = 6.69e+04 \cdot \exp\left( -\frac{(f - 224609.375)^2}{2 \cdot 35.86^2} \right) \cdot e^{-i \cdot 3.142} + 3.52e+04
\end{equation}
Спектральная компонента сигнатуры $\Psi_{27}$ ($f = 251953.125$ Гц):
\begin{equation}
S(f) = 3.81e+04 \cdot \exp\left( -\frac{(f - 251953.125)^2}{2 \cdot 74.45^2} \right) \cdot e^{-i \cdot 3.142} + 2.80e+04
\end{equation}
Спектральная компонента сигнатуры $\Psi_{28}$ ($f = 300781.250$ Гц):
\begin{equation}
S(f) = 1.83e+04 \cdot \exp\left( -\frac{(f - 300781.250)^2}{2 \cdot 157.34^2} \right) \cdot e^{-i \cdot 0.000} + 1.58e+04
\end{equation}
Спектральная компонента сигнатуры $\Psi_{29}$ ($f = 390625.000$ Гц):
\begin{equation}
S(f) = 1.25e+05 \cdot \exp\left( -\frac{(f - 390625.000)^2}{2 \cdot 181.40^2} \right) \cdot e^{-i \cdot 3.142} + 1.10e+05
\end{equation}
Спектральная компонента сигнатуры $\Psi_{30}$ ($f = 8203125.000$ Гц):
\begin{equation}
S(f) = 3.05e+04 \cdot \exp\left( -\frac{(f - 8203125.000)^2}{2 \cdot 69.28^2} \right) \cdot e^{-i \cdot 3.142} + 2.18e+04
\end{equation}
Спектральная компонента сигнатуры $\Psi_{31}$ ($f = 15625000.000$ Гц):
\begin{equation}
S(f) = 9.44e+03 \cdot \exp\left( -\frac{(f - 15625000.000)^2}{2 \cdot 52.49^2} \right) \cdot e^{-i \cdot 3.142} + 6.09e+03
\end{equation}
Спектральная компонента сигнатуры $\Psi_{32}$ ($f = 22265625.000$ Гц):
\begin{equation}
S(f) = 5.09e+03 \cdot \exp\left( -\frac{(f - 22265625.000)^2}{2 \cdot 85.71^2} \right) \cdot e^{-i \cdot 3.142} + 3.89e+03
\end{equation}
Спектральная компонента сигнатуры $\Psi_{33}$ ($f = 28906250.000$ Гц):
\begin{equation}
S(f) = 1.58e+03 \cdot \exp\left( -\frac{(f - 28906250.000)^2}{2 \cdot 143.30^2} \right) \cdot e^{-i \cdot 3.142} + 1.34e+03
\end{equation}
Спектральная компонента сигнатуры $\Psi_{34}$ ($f = 36718750.000$ Гц):
\begin{equation}
S(f) = 7.44e+02 \cdot \exp\left( -\frac{(f - 36718750.000)^2}{2 \cdot 237.12^2} \right) \cdot e^{-i \cdot 3.142} + 6.75e+02
\end{equation}
Спектральная компонента сигнатуры $\Psi_{35}$ ($f = 47265625.000$ Гц):
\begin{equation}
S(f) = 3.85e+02 \cdot \exp\left( -\frac{(f - 47265625.000)^2}{2 \cdot 1663.20^2} \right) \cdot e^{-i \cdot 3.142} + 3.80e+02
\end{equation}
Спектральная компонента сигнатуры $\Psi_{36}$ ($f = 71484375.000$ Гц):
\begin{equation}
S(f) = 2.37e+02 \cdot \exp\left( -\frac{(f - 71484375.000)^2}{2 \cdot 68.12^2} \right) \cdot e^{-i \cdot 3.142} + 1.69e+02
\end{equation}
Спектральная компонента сигнатуры $\Psi_{37}$ ($f = 82421875.000$ Гц):
\begin{equation}
S(f) = 1.82e+02 \cdot \exp\left( -\frac{(f - 82421875.000)^2}{2 \cdot 64.66^2} \right) \cdot e^{-i \cdot 3.142} + 1.27e+02
\end{equation}
Спектральная компонента сигнатуры $\Psi_{38}$ ($f = 89843750.000$ Гц):
\begin{equation}
S(f) = 1.67e+02 \cdot \exp\left( -\frac{(f - 89843750.000)^2}{2 \cdot 35.86^2} \right) \cdot e^{-i \cdot 3.142} + 8.80e+01
\end{equation}
Спектральная компонента сигнатуры $\Psi_{39}$ ($f = 100781250.000$ Гц):
\begin{equation}
S(f) = 9.53e+01 \cdot \exp\left( -\frac{(f - 100781250.000)^2}{2 \cdot 74.45^2} \right) \cdot e^{-i \cdot 3.142} + 6.99e+01
\end{equation}
Спектральная компонента сигнатуры $\Psi_{40}$ ($f = 120312500.000$ Гц):
\begin{equation}
S(f) = 4.56e+01 \cdot \exp\left( -\frac{(f - 120312500.000)^2}{2 \cdot 157.34^2} \right) \cdot e^{-i \cdot 0.000} + 3.94e+01
\end{equation}
\end{document}
